% Generated by `taskgen course examples` from lattices.gauss (seed 8).
% Do not edit: change examples.yaml and regenerate.
\begin{example}[алгоритм Гаусса]\label{ex:lattice-gauss}
В решётке $\mathcal{L}$, заданной базисными векторами \[v_1 = \begin{pmatrix}39 \\ -39\end{pmatrix}, v_2 = \begin{pmatrix}68 \\ -65\end{pmatrix},\] найти кратчайший вектор методом Гаусса.

\par\noindent\textit{Решение.} $\det \mathcal{L} = 117$, $HR(\mathcal{L}) = \left(\frac{|\det\mathcal{L}|}{\|v_1\|\|v_2\|}\right)^{1/2} = 0.150$ --- базис плохой. \par Алгоритм Гаусса: если $\|v_2\| < \|v_1\|$ "--- меняем местами (${}^{\leftrightarrow}$); $m = \lfloor v_1 \cdot v_2 / \|v_1\|^2 \rceil$; если $m = 0$ "--- стоп, иначе $v_2 \gets v_2 - m v_1$.
\[\begin{array}{c|r|r|r|r|r}
 & v_1 & v_2 & \|v_1\|^2 & v_1\cdot v_2 & m \\ \hline
1 & (39, -39) & (68, -65) & 3042 & 5187 & 2 \\
2^{\leftrightarrow} & (-10, 13) & (39, -39) & 269 & -897 & -3 \\
3^{\leftrightarrow} & (9, 0) & (-10, 13) & 81 & -90 & -1 \\
4 & (9, 0) & (-1, 13) & 81 & -9 & 0 \\
\end{array}\]
Кратчайший вектор --- $v_1$ последней строки, $\|v_1\|^2 = 81$.\[r = \pm(9, 0)^T,\quad HR(\mathcal{L}) = 0.150.\]
\end{example}
